% gtwillman.tex -- the GTWILL manual.  Process with luatex, which can find
% gwebmac.tex and hangulfont.tex, for example
%       TEXINPUTS=.:gweb: luatex gtwillman.tex
% or simply `make manual' from the top of the repository.
%
% It is formatted like GWEB's gwebman.tex, which is formatted in the manner of
% Knuth and Levy's cwebman.tex: it loads the GWEB macro package and adds its own
% section, list, verbatim, and running-head macros on top. The Korean
% translation, gtwillman-ko.tex, keeps the same macros; change it with this one.
% Only GWEB's hangulfont.tex is loaded, not kotexgweb.tex, whose baselines
% are for Korean text: the code displays pair \tentt with its Hangul
% typewriter \tentz, for the line of Korean in `Weaving in Korean'.

\input gwebmac
\input hangulfont
\pdffalse \dvipdffalse        % a plain manual: no hyperlinks or bookmarks
\parskip=0pt plus 1pt
\parindent=18pt
\abovedisplayskip=.5\abovedisplayskip  \belowdisplayskip=.5\belowdisplayskip
\abovedisplayshortskip=.5\abovedisplayshortskip
\belowdisplayshortskip=.5\belowdisplayshortskip

\font\eighttt=cmtt8
\font\ninett=cmtt9
\font\twelvett=cmtt12
\font\authorfont=cmr12
\font\sectionfont=cmbx12
\font\eightit=cmti8
\font\eightbf=cmbx8
\font\logosl=logosl10 % the METAFONT logo, slanted for a book title
\def\slMF{{\logosl META}\-{\logosl FONT}}

%---inline code; @ < > = etc. are safe in \tt, so avoid braces inside \cw ---
\def\cw#1{{\tt #1}}
\def\bs{\char92 }                              % a literal backslash, in \tt
\def\lb{\char123 }\def\rb{\char125 }           % literal braces, in \tt

%---a centred chapter heading, in the manner of cwebman's \section ---
\def\chapter#1{\penalty-500\bigskip
  \centerline{\sectionfont\def\.##1{{\twelvett##1}}#1}\nobreak\vskip 6pt}
%---a run-in bold subsection heading ---
\def\section#1{\medskip\noindent{\bf #1.}\enspace\ignorespaces}

%---a bulleted list item ---
\def\bull{\smallskip\item{$\bullet$}}

%---a verbatim code display (TeXbook recipe): | is the escape inside it.
%   Each line is a paragraph of its own; a \nobreak after each keeps the
%   display on one page ---
\def\begincode{\par\medskip\begingroup
  \def\do##1{\catcode`##1=12 }\dospecials
  \catcode`\|=0 \def\par{\endgraf\nobreak}%
  \obeylines \obeyspaces \tentt\tentz \parindent=0pt \leftskip=\parindent\relax}
\def\endcode{\endgroup\medskip\noindent}

%---a sample mini-index: three columns under a rule, as gtwimac sets one ---
\def\sampleentry#1{\hbox{#1}}
\def\samplemini#1#2#3{\par\medskip
  \hrule\kern4pt
  \hbox to\hsize{\eightrm \def\it{\eightit}\def\bf{\eightbf}\let\tt=\eighttt
    \baselineskip=9pt \vtop{#1}\hfil\vtop{#2}\hfil\vtop{#3}}%
  \medskip\noindent}

%---running heads, in the manner of cwebman ---
\def\runninghead{{\ninett GTWILL} USER MANUAL (VERSION 0.10.7)}
\def\lheader{\mainfont\folio\hfil{\sc\runninghead}\hfil}
\def\rheader{\hfil{\sc\runninghead}\hfil\mainfont\folio}

% =============================================================================
% Title page (no running head), in the manner of cwebman.
\headline={\hfil}\footline={\hfil}
\null\vfill
\centerline{\titlefont The {\ttitlefont GTWILL} System}
\vskip 6pt
\centerline{\sectionfont Mini-indexes for Literate Go Programs}
\vskip 18pt
\centerline{(Version 0.10.7)}
\vskip 24pt
\centerline{\authorfont modeled on Knuth's {\ttitlefont CTWILL}}
\vfill
\noindent
This is a guide to \.{GTWILL} for readers who already know \.{GWEB}. \.{GTWILL}
weaves a \cw{.w} file exactly as \cw{gweave} does, but it sets the pages in
facing pairs and puts a {\it mini-index\/} at the foot of every right-hand page:
the identifiers that the two pages use but do not define, each with a word about
what it is and the section that defines it. It is to \.{GWEB} what Knuth's
\.{CTWILL} is to \.{CWEB}, and \.{CTWILL} descends from \.{TWILL}, the private
\.{WEB} program with which Knuth produced {\sl \TeX: The Program\/} and {\sl
\slMF: The Program\/} in 1985. Knuth tells the story in
``Mini-indexes for literate programs,'' reprinted in {\sl Digital
Typography\/} (1999), 225--245. A web needs no change to be woven by
\cw{gtwill}; but good mini-indexes, as Knuth found, take some work, and most of
this manual is about that work.
\pageno=0
\vfil\eject

% Running heads from the second page on.
\headline={\ifodd\pageno\rheader\else\lheader\fi}

\chapter{\.{GTWILL} at a glance}
\noindent\.{GTWILL} mirrors \.{CTWILL} tool for tool:

\bull \cw{gtwill} $\approx$ \cw{ctwill}---weaves a web, with mini-index
information;

\bull \cw{gtwill-refsort} $\approx$ \cw{refsort}---sorts the mini-index entries
between two \TeX\ passes;

\bull \cw{gtwill-twinx} $\approx$ \cw{twinx}---merges the indexes of several
programs into one master index;

\bull \cw{gtwimac.tex} $\approx$ \cw{ctwimac.tex}---the page builder for facing
pages with mini-indexes;

\bull \cw{gtproofmac.tex} $\approx$ \cw{proofmac.tex}---proofmode, in which each
section is followed by its entries;

\bull \cw{gtwinxmac.tex} $\approx$ \cw{twinxmac.tex}---the master index.

\noindent The whole cycle for one web \cw{foo.w}, with its change file
\cw{foo.ch}, is
\begincode
gtwill foo.w foo.ch                   # -> foo.tex, foo.aux
luatex foo                            # first pass: -> foo.pdf, foo.mref
gtwill-refsort < foo.mref > foo.smref
luatex foo                            # second pass: mini-indexes alphabetized
|endcode
and \cw{pdftex} serves as well as \cw{luatex}, unless the web is in Korean.

\chapter{What a mini-index says}
\noindent Here is the foot of a right-hand page of {\sl \TeX: The Program},
translated into \GO/ and set by \.{GTWILL}:
\samplemini
 {\sampleentry{{\it appendChar\/}: {\bf func}( ), \S53.}
  \sampleentry{{\it curLength\/}: {\bf func}( ) {\bf int32}, \S52.}
  \sampleentry{{\it flushChar\/}: {\bf func}( ), \S53.}
  \sampleentry{{\it hh\/}: {\bf twoHalves}, \S129.}}
 {\sampleentry{{\it mem\/}: [\,]{\bf memoryWord}, \S132.}
  \sampleentry{{\it memEnd\/}: {\bf pointer}, \S134.}
  \sampleentry{{\it memMin\/} = 0, \S13.}
  \sampleentry{{\bf memoryWord} = {\bf struct}, \S129.}}
 {\sampleentry{{\it printCurrentString\/}: {\bf func}( ), \S85.}
  \sampleentry{{\it printLn\/}: {\bf func}( ), \S69.}
  \sampleentry{{\it printStr\/}: {\bf func}( ), \S72.}
  \sampleentry{{\it showBox\/}: {\bf func}( ), \S218.}}%
Each entry is an identifier that the two facing pages use but do not define,
set the way the index sets it (bold for a type, typewriter for an \cw{@d} name
or a constant in capitals, italic otherwise); then its {\it meaning}, after a
colon, or after an equals sign for a constant or a type; then the section that
defines it. A name defined in another program cites that program's title before
the section number, {\sc GBGRAPH}\thinspace\S12, and a name from a package or a
file cites it in typewriter, \cw{fmt}.

\section{Current meanings}
\cw{gtwill} works the way \.{CTWILL} does. Every identifier has a current
meaning, consisting of its abstract type and the number of the section in which
it was most recently defined. If your program says \cw{var s []byte} in
section~3, the meaning of \cw{s} becomes ``\cw{[]byte}, \S3'' while \cw{gtwill}
is processing that section. If you refer to \cw{s} in section~10, and \cw{s}
hasn't been redefined in the meantime, and section~10 doesn't wind up on the
same spread as section~3, the mini-index of section~10's spread will say
``{\it s\/}: [\,]{\bf byte}, \S3.'' At the end of each section \cw{gtwill} gives
\TeX\ the list of identifiers the section uses, with their meanings, and the
macros of \cw{gtwimac.tex} decide which of them belong in each mini-index, based
on how the pages break.

\section{Meanings from \GO/ declarations}
\CEE/ spells a declaration's type for all to see; \GO/ often leaves it to the
compiler, and \cw{gtwill} is no compiler. \.{CTWILL} reads a definition off the
partly reduced scraps of \cw{cweave}'s grammar; \cw{gweave} has no grammar, only
the tokens around a defined name, so \cw{gtwill} reads the declaration directly:

\bull a \cw{var}, a field, or a parameter means its type;

\bull a \cw{type} means $=$ and its definition, with the body of a \cw{struct}
or \cw{interface} elided;

\bull a \cw{const} means $=$ and its value, when that is a literal or a short
expression (\cw{memTop - 11}), and its type or just \cw{const} otherwise;

\bull a function means \cw{func( )} and its results, the parameters elided; a
method means \cw{(*T) func( )} and its results; a label means \cw{label};

\bull a short declaration such as \cw{n := len(s)} states no type, so \cw{gtwill}
guesses one from the form of the value---a literal, a composite literal, a call
of \cw{make} or \cw{new}, a conversion, a function literal---and otherwise
says \cw{var}.

\noindent Parameter lists and long array lengths are elided as \.{CTWILL}
elides them. A field is recognized as one inside a \cw{struct}'s body, also
after a section reference there, and in a section that a \cw{struct} splices in
among its fields.

\section{What is not cited}
An identifier that nobody has defined is cited as `???', so that you notice it,
but only if it is an ordinary name of the program. \cw{gtwill} keeps quiet about
predeclared names (\cw{int}, \cw{len}), the names of packages, imported or the
program's own, and names selected with a dot, like the \cw{Println} of
\cw{fmt.Println} or the \cw{Len} of \cw{b.Len()}, which belong to some other
package. Give such a name a meaning with \cw{@\$} (see below) and it is cited
like any other.

\chapter{Running the programs}
\begincode
gtwill [-P] [-o dir] file[.w] [change[.ch]]
gtwill-refsort < file.mref > file.smref
gtwill-twinx file.tex ... > index.tex
|endcode
\cw{gtwill} takes the arguments and flags of \cw{gweave}---the \cw{.w} and
\cw{.ch} may be omitted, \cw{-o} chooses the output directory, \cw{-version}
prints the version---and one more, \cw{-P}, for proofmode (see below). Besides
\cw{foo.tex} it writes \cw{foo.aux}, which records the web's own definitions
in a form that other webs can read. The woven file begins with \cw{\bs input
gtwimac}, which loads \cw{gwebmac}; a stray \cw{\bs input gwebmac} in the limbo is
dropped.

\section{One run, two passes}
\.{CTWILL} had to be run twice, ``once to prime the pump and once to get decent
answers.'' \cw{gweave} already makes a silent pass over the whole web before the
pass that writes its output, and \cw{gtwill} primes the pump in that first pass,
so a single run gives the decent answers.

\TeX\ must still run twice. The first pass says `This is the first pass!' and
sets pages that look final except that the entries of each mini-index come in
the order the sections cited them; it also writes those entries to
\cw{foo.mref}. \cw{gtwill-refsort} sorts them, page by page, into
\cw{foo.smref}, and the second pass, finding that file, says `This is the second
pass!' and sets them alphabetized. (\.{CTWILL} calls these files \cw{.ref} and
\cw{.sref}; \.{GTWILL} adds an \cw{m}, because a web is apt to keep a \cw{.ref}
file of its own.) Should you change the web, delete \cw{foo.smref} and start
over: it describes pages that no longer exist. If it ends early, \TeX\ says so in
the log. Saying \cw{\bs let\bs norefs=\bs relax} before \cw{\bs input gtwimac}, that
is, on the command line, reads and writes neither file:
\begincode
luatex "\let\norefs=\relax \input foo"
|endcode

\section{Proofmode}
\cw{gtwill -P} makes the woven file load \cw{gtproofmac.tex} instead of
\cw{gtwimac.tex}. \TeX\ then sets \cw{gweave}'s ordinary pages, without
mini-indexes, and follows each section with the entries it generates. Check
your hints there before making mini-indexes: mini-indexes can mask errors when
the page breaks are favorable, and the errors may reveal themselves later, after
your program has changed. The proofmode output is much easier to check than the
mini-indexes themselves.

\chapter{Hints in the change file}
\noindent The meanings \cw{gtwill} concocts are not always right. Your
exposition may talk about \cw{f(x)} without meaning program variables \cw{f} and
\cw{x}; a name may be defined in another web, unknown to \cw{gtwill}; a name may
be defined in several places, and \cw{gtwill} won't know which definition to
choose. Book production with \.{CTWILL} always involved a change file of hints,
and so does it here. These control codes are \cw{gtwill}'s alone; \cw{gtangle}
skips them, which is no problem, since you use them only in the change files you
set up for publishing.

\section{\cw{@\$}: a meaning}
The operation
\begincode
@$s {FOO}3 \KW{int}@>
|endcode
changes the current meaning of \cw{s} to the \TeX\ text \cw{\bs KW\lb int\rb},
defined in section~3 of the program whose title is {\sc FOO}. It has four
components: an identifier, followed by a space; a program title, in braces; a
section number, followed by a space; and a \TeX\ part. If the \TeX\ part starts
with \cw{=}, the entry has an equals sign instead of a colon:
\cw{@\$bufSize \lb PROG\rb 10 =\bs NU\lb 200\rb @>} gives ``{\it bufSize\/} $=$
200, \S10'' or ``{\it bufSize\/} $=$ 200, {\sc PROG}\thinspace\S10,'' depending on
whether {\sc PROG} is the title of the current program. If the \TeX\ part is
\cw{\bs zip}, the entry has neither colon nor equals sign, just a comma. The
title of the current program is what its limbo gives to \cw{\bs def\bs title};
failing that, it is the base name of the output file.

The title and section number can also be replaced by a string in doublequotes:
\cw{@\$Println "fmt" \bs zip@>} gives ``{\it Println}, \cw{fmt}.'' Such a string is
set in typewriter as it stands; you need not escape the characters special to
\TeX\ in it, since \cw{gtwill} escapes them, so that a name from the file
\cw{"gb\_graph.w"} can say so.

\section{\cw{@\%}: temporary meanings}
After \cw{@\%}, the \cw{@\$} operations up to the beginning of the next section
define {\it temporary\/} meanings instead of permanent ones. A temporary meaning
goes into the mini-index of the current section; the permanent meaning of the
identifier is not changed, nor does it appear in that section's entries. If
several temporary meanings are given to one identifier in a section, all of them
appear. Each \cw{@\%} toggles the convention, so after an even number of them
the meanings are permanent again. (In \cw{gweave} and \cw{gtangle}, \cw{@\%}
begins a comment that runs to the end of the line; \cw{gtwill} gives it
\.{CTWILL}'s meaning instead, so a web woven by \cw{gtwill} should not use it as a
comment.)

\section{\cw{@-}: no entry}
\cw{@-foo@>} says that \cw{foo} should not generate a mini-index entry in the
current section, unless a temporary meaning is assigned. If a later
declaration in that section would give \cw{foo} a new permanent meaning, the
current meaning is not redefined, and that current meaning, which an
\cw{@\$foo ...@>} may have set, is what goes into the \cw{.aux} file. So you can
keep \cw{gtwill} from contaminating the \cw{.aux} file with things you don't
like.

\section{\cw{@r}: a right-hand page}
\cw{@r} (or \cw{@R}) puts the macro \cw{\bs shortpage} before the next section,
which then appears at the top of a right-hand page, if it would ordinarily have
appeared near the bottom of a left-hand page and split across the pages. If the
next section is starred, \cw{@r} inverts the left/right logic: the starred
section is postponed to a left-hand page if it would begin on a right-hand page,
or vice versa. Both are fragile; use them only when fine-tuning a set of pages.
The macro \cw{\bs forcebr}, in a \TeX\ part, ends the spread before the next
section.

\section{The \cw{.aux} and \cw{.bux} files}
The current meaning of every identifier is initially \cw{\bs uninitialized}.
Then \cw{gtwill} reads the file \cw{system.bux}, if there is one; then the
definitions the web itself makes; then \cw{foo.bux}. The \cw{.bux} files are
looked for in the output directory, in the web's own directory, and in the
current directory, in that order. Like the \cw{.aux} files, they should contain
only \cw{@\$} specifications, as the name says: `bux' is the plural of `\$'. They
may also contain \cw{@i} includes.

Every entry of \cw{foo.aux} is such an \cw{@\$} operation, one per line, so you
can paste entries from it into a change file whenever you want to tell
\cw{gtwill} about definitions that are out of order, and into the \cw{.bux} file
of another web that uses them: that is how the programs of a book refer to one
another. The convention is to put the definitions of library functions into
\cw{system.bux}, and those of the specifically foo-ish packages that \cw{foo.w}
imports into \cw{foo.bux}. For a book of several programs, run \cw{gtwill} on
every web once to have all the \cw{.aux} files, make the \cw{.bux} files from
them, and then weave and typeset each web in earnest.

\chapter{Keeping the hints outside the web}
\noindent Hints written by hand into a change file soon become a nuisance in
two ways. One is the section numbers: those written in an \cw{@\$} go wrong
whenever a section is added to the web or taken out. The other is the \TeX: a
meaning written by hand, like \cw{\bs KW\lb int\rb}, easily drifts from the
form in which \cw{gtwill} sets the meanings it reads off declarations. So the
\GO/ translations of Knuth's programs \cw{go-sgb}, \cw{go-mmix}, and
\cw{go-advent} keep a directory \cw{twill} in the repository, and everything
the mini-indexes need lives there. The original is \cw{go-sgb}'s; \cw{go-mmix}
adapted it, and \cw{go-advent} cut it down for a single web. This chapter takes
\cw{go-advent} as its example.
\begincode
twill/hints.txt     # the hints, one name per line, kept by hand
twill/bux.py        # writes system.bux and advent.bux before gtwill runs
|endcode
These two files are all that is committed. The two \cw{.bux} files that
\cw{bux.py} makes are generated, so they go in \cw{.gitignore}, and to change
them you change \cw{hints.txt}. The web itself never changes for the sake of a
hint, so what \cw{gtangle} and \cw{gweave} see stays as it is.

\section{The form of a hint}
Each line of \cw{hints.txt} gives a name, a place, and a meaning. Here are two
of \cw{go-advent}'s hints: names of Knuth's \.{CWEB} original that the
commentary mentions in \cw{|...|} but the \GO/ code does not have, given as
their meaning the \GO/ that takes their place.
\begincode
change_to "advent.w" =changeTo
fgets "C" -
|endcode
The place is written thus:

\bull \cw{\S}---the section of this web in which the name first appears inside
\cw{|...|}. The script finds the number every time the pages are made, so adding
or removing sections does not throw it off.

\bull \cw{\S12}---section~12. The script checks that the name really appears in
that section's \cw{|...|}.

\bull \cw{gbwords\S}---a section of another web, as used in \cw{go-sgb} and
\cw{go-mmix}, which have many.

\bull \cw{"advent.w"}---a doublequoted string: the name of an original file or a
\GO/ package, and \cw{"C"} for the \CEE/ language or its standard library.

\noindent The meaning is written in \GO/, not in \TeX. With an equals sign
(\cw{=}) in front it is set after an equals sign, with a colon (\cw{:}) after a
colon, and a hyphen (\cw{-}) alone sets only the place (\cw{\bs zip}). The script
weaves these pieces of \GO/ with \cw{gweave}, together with the web's limbo, and
puts the resulting \cw{\bs PB\lb...\rb} into the \cw{@\$} operation. So a meaning
given as a hint is set exactly like one read off a declaration, and the types the
web defines come out bold there too.

\section{The two \cw{.bux} files}
Before \cw{gtwill} runs, \cw{bux.py} writes two \cw{.bux} files. The file
\cw{advent.bux} gets the hints of \cw{hints.txt}. The file \cw{system.bux} gets
the names that the web uses without defining them: in \cw{go-advent}, the
variables, with their types, declared by \cw{cavedata.go} and \cw{objdata.go},
data dumped from the \CEE/ original. Those files do not come from a web, so their
names would otherwise be cited as `???'. With several webs \cw{system.bux} has
more to do. For each package that a web imports, \cw{go-mmix}'s script collects
from the \cw{.aux} files of that package's webs only the exported names, since
those files hold local variables too. It also collects the names declared by
sibling webs of the same package, and drops altogether a name that two places
define differently: better not to cite a name than to cite it wrong.

\cw{gtwill} reads \cw{system.bux}, the web's own definitions, and
\cw{advent.bux}, in that order. So the web's own definitions win over imported
names, and the hints win over both. The two files are written where \cw{gtwill}
looks for them, in the web's own directory.

\section{\cw{make twill}}
\cw{go-advent}'s \cw{Makefile} wraps it all up in \cw{make twill}:
\begincode
python3 twill/bux.py advent              # -> system.bux, advent.bux
gtwill advent.w                          # -> advent.tex, advent.aux
luatex --jobname=advent-twill advent.tex # first pass
gtwill-refsort < advent-twill.mref > advent-twill.smref
luatex --jobname=advent-twill advent.tex # second pass
|endcode
\cw{gtwill} writes \cw{advent.tex}, as \cw{gweave} does; typesetting it under
the job name \cw{advent-twill} keeps the \cw{.pdf}, \cw{.log}, \cw{.mref}, and
\cw{.smref} apart from those of \cw{make doc}. The old \cw{.smref} and \cw{.pdf}
are removed before the first pass: the old \cw{.smref} describes pages that no
longer exist, and an old \cw{.pdf} would hide a failed run. At the end two
numbers are counted: the typesetting warnings in the log, and the
\cw{\bs uninitialized} entries left in the \cw{.smref}, those set as `???' for
want of a meaning. The work is done when both are zero. With several webs, as the
previous chapter said, first run \cw{gtwill} on every web once to have all the
\cw{.aux} files, and then run the script, \cw{gtwill}, and the two passes on each
web in turn.

\chapter{The pages}
\noindent \cw{gtwimac.tex} keeps everything of \cw{gwebmac.tex}---fonts,
displayed code, links, index, contents---except the page builder. Each section is
set by itself and added to the current spread unless the spread would then
overflow; in that case the spread is shipped first, as a left-hand page cut off
by \cw{\bs vsplit} and a right-hand page holding the rest and the mini-index. The
left-hand page's running head shows the page number, the current starred
section's title, the program's title, and the first section of the spread; the
right-hand page's shows the last section, the title, and the page number.

\section{Where it differs from \.{CTWILL}}
\bull The page builder measures the mini-index instead of estimating it from the
number of entries, and tries the left-hand page's \cw{\bs vsplit} before it
accepts a section, so the right-hand page never overflows. A section longer than
two pages continues on the next spread instead of overflowing, and there is no
limit on the number of sections in a spread.

\bull The index section, the index, the list of section names, and the contents
are set as \cw{gweave} sets them: \cw{gtwill} writes \cw{\bs twilllast} before the
last starred section, and \cw{gtwimac.tex} ships the last spread and goes back to
\cw{gwebmac}'s ordinary pages. (\.{CTWILL} sets no index at all; Knuth gathered
the indexes of a book's programs into one with \cw{twinx}.)

\bull A change file meant for \cw{gtwill} is a big one, touching nearly
everything, so the sections it touches are not starred, nor listed before the
index, as \cw{gweave} would star and list them; and since none is marked as
changed, a \cw{\bs let\bs maybe=\bs iffalse} in the limbo hides nothing.

\chapter{Tuning}
\noindent The pages are \cw{gwebmac}'s \cw{\bs pagewidth} by
\cw{\bs pageheight}. These parameters, set in the limbo, shape the rest:

\bull \cw{\bs miniextra} (1pc)---how much lower the right-hand page reaches, to
make room for its mini-index;

\bull \cw{\bs minisep} (11pt) and \cw{\bs minirulesep} (4pt)---the space above and
below the rule over a mini-index;

\bull \cw{\bs minicolsep} (1pc)---the space between its three columns;

\bull \cw{\bs minibase} (9pt), \cw{\bs minitop} (7pt), and \cw{\bs minidepth}
(3pt)---its baselineskip, its topskip, and the room below its last line;

\bull \cw{\bs leftstretch}---how readily a left-hand page ends early at a section
boundary. A page of code has no stretch, and without some every break would be
equally bad.

\noindent The words of an entry are macros too, for a localization to
redefine: \cw{\bs minicolon} (`: '), \cw{\bs miniequals} (`\thinspace$=$\thinspace'),
\cw{\bs twilllabel} (`label'), and \cw{\bs uninitialized} (`???'). A decoration
drawn behind every page goes in \cw{\bs twillbackground}, which comes first on
each page of a spread, as it would in a \cw{\bs plainoutput} of your own; it must
take no room.

As \cw{ctwimac.tex} does, \TeX\ writes to its log, for every section, how much
room the spread would need with that section in it, and
`(shortpage)' where a \cw{\bs shortpage} took effect. A section that runs on to
the next spread is reported too. With these numbers and \cw{@r} and
\cw{\bs forcebr}, a set of pages can be fine-tuned.

\chapter{The master index}
\begincode
gtwill-twinx gbgraph.tex gbbasic.tex ... > index.tex
luatex index
|endcode
\cw{gtwill-twinx} reads the index that \cw{gtwill} (or \cw{gweave}) puts at the end
of each \cw{.tex} file, and the title from its \cw{\bs def\bs title\lb...\rb},
and merges them: each entry cites the sections of each program under its title,
in the order the files were given, with the sections of definition underlined.
The result loads \cw{gtwinxmac.tex}, which sets it in two columns with
\cw{gwebmac}'s index macros. If a file \cw{gtwinx-startup.tex} can be found,
\cw{gtwinxmac.tex} reads it: that is the place for \cw{\bs input kotexgweb}, for a
\cw{\bs def\bs title} (the default is `Index'), for a new \cw{\bs twinxtop} (the
heading) or \cw{\bs twinxnext} (what goes between the citations of two programs),
and for the macros that the programs' limbos define and their index entries use.

\chapter{Weaving in Korean}
\noindent A Korean web is woven by \cw{gtwill} as by \cw{gweave}: put
\cw{\bs input kotexgweb} in its limbo and typeset with \cw{luatex}. The woven file
reads \cw{gtwimac.tex} first and \cw{kotexgweb.tex} afterwards, and
\cw{kotexgweb}'s own \cw{\bs eightpoint}, with its Hangul faces, replaces
\cw{gtwimac}'s for the mini-indexes. The word \cw{\bs twilllabel} stays
English unless you redefine it, in the limbo after the \cw{\bs input}:
\begincode
\input kotexgweb
\def\twilllabel{레이블}
|endcode

\chapter{Installing \.{GTWILL}}
\begincode
make            # tangle the webs, build bin/gtwill and the two utilities
make test       # go vet + go test
make install    # the commands into GOBIN, the TeX macros into TEXMFHOME
|endcode
Only \GO/ is needed to build: \cw{gtangle} comes from the \.{GWEB} module pinned in
\cw{go.mod} (\cw{go tool gtangle}), and the master \cw{gweave.w} is reached through
a symbolic link \cw{gweb} into the module cache, which \cw{make} creates. To
typeset you need \.{GWEB}'s \cw{gwebmac.tex} on the \TeX\ path as well
(\.{GWEB}'s \cw{install.sh} puts it there), and a \TeX\ engine with e-\TeX:
\cw{luatex} or \cw{pdftex}.

\chapter{\.{GTWILL} is written in itself}
\noindent Like \.{CTWILL}, which is \cw{cweave} with a change file
(\cw{ctwill.ch}), \cw{gtwill} is \cw{gweave} with a change file: \cw{gtwill.ch}
applies to \cw{cmd/gweave/gweave.w} of the pinned \.{GWEB} release, and
\cw{gtangle} makes \cw{gtwill.go} of the two. The two utilities are webs of their
own, \cw{refsort.w} and \cw{twinx.w}. The tests live in the webs too.
\begincode
make doc        # the listings, woven by gweave, into build/
make selfdoc    # gtwill's own listing with mini-indexes, into build/self/
make manual     # this manual; make manual-ko for the Korean one
|endcode
\cw{make doc} weaves \cw{gweave.w} with \cw{gtwill.ch} the ordinary way, with the
changed sections starred; its second section is a manual in Knuth's words for
\.{CTWILL}, kept where they still fit. \cw{make selfdoc} sets the same listing
with \cw{gtwill} itself, as Knuth set \.{CWEB}'s with \.{CTWILL}.

\vfill
\centerline{\eightrm {\eighttt GTWILL}---mini-indexes for literate Go programs.}
\bye
